% !TEX root = handbook-physics.tex
% !TEX TS-program = pdflatex
% !TEX encoding = UTF-8 Unicode




\chapter[Structure Of Solids]{Structure and Certain Properties of Solids}
\label{ch-40}

\section{Structure of Solids}
\label{sec-40.1} 

40.1.1 Solids are bodies having constant shape and volume.
Distinction is made between crystalline and amorphous solids.
Externally, crystals have a regular geometric shape. They have
a periodically repeated arrangement of their constituent particles
throughout the whole crystal. This arrangement is called the
crystal, or space, lattice. In this sense, a crystal is said to have
long-range order (cf.- Sect. 13.1.2). Crystals are bounded by flat
faces, orderly arranged with respect to one another, which converge
at the edges and vertices. Large separate crystals having the
shape of regular polyhedrons are called monocrystals. Their shape
is determined by their chemical composition. Poly crystals have
a fine crystalline structure made up of a great number of small
randomly arranged crystals (crystal grains or crystallites) that have
grown together.
Amorphous solids (pitch, glass, etc.) are supercooled liquids and
have no clearly defined properties of crystals.
In the following, discussion is confined to crystalline bodies.
40.1.2 Each particle in a crystal lattice is subject to the forces
of intermolecular interaction (12.1.4). An equilibrium arrangement
of all the particles of a solid at the points of.the crystal lattice
(8.1.4) corresponds to minimum free energy of the crystal (11.4.5)
and its most stable state. The particles at the points of the crystal
laltice are located at certain equilibrium distances from one another
called the lattice constants.
40.1.3 The main types of crystalline solids, distinguished from
one another by the nature of the forces of interaction between
their particles and the kinds of particles located at the crystal
lattice points, are the following.
37*mo Structure aNd pftoPEMiES ob SoLtbs [/,o
(a) Ionic crystals (NaCl, calcium carbonate and other salts).
Located at the lattice points are regularly alternating positive
and negative ions, between which are heteropolar bonds (39.4.4).
(b) Valence (covalent) crystals (G, Ge, Te, etc.). The neutral atoms
at the lattice points have homopolar bonds (39.4.5) between them.
Such crystals are found in semiconductors (41.10.1) and in many
organic solids.
(c) Molecular crystals (At, CH4, paraffin, etc.). The points of the
crystal lattice are occupied by molecules that are bound by Van
der Waals forces, mainly of the dispersion type (12.1.8).
(d) Metals (Na, Cu, Al, etc.). At the lattice points arc positive
ions that are formed after the outer (valence) electrons (39.3.9)
have been stripped from the atoms. These electrons form an
electron gas (16.4.1) of collectivized free particles.
The metallic bond is a specific type of chemical bond (39.4.1)
that is developed between the ions in the crystal lattice and the
electron gas. The electrons “constrict” the positive ions (mainly
by electrostatic forces) and thereby compensate for repulsion
between the ions. At distances between the ions equal to the
lattice constant (40.1.2), the metallic crystal acquires a stable
state.
40.1.4 The anharmonic nature of the thermal vibration of the
particles at the crystal lattice points (8.1.4) is manifested in the
fact that the dependence of the potential energy of interaction of the
particles on their displacement from the equilibrium position is not
a parabolic one. The force acting on a particle is not quasi-elastic
(40.3.5). These facts are of prime importance in understanding
certain thermal (40.2.4) and electrical (41.5.4) properties of solids.
40.1.5 A typical property of monocrystals (40.1.1) is their anisotropy
(crystalline anisotropy), which is the dependence of the physical
properties of solids (thermal, elastic, electrical and optical
properties) on the direction in the crystal.

\section{Thermal Expansion of Solids}
\label{sec-40.2} 

40.2.1 A solid is subject to thermal expansion, of either linear or
volume type, when its temperature is raised. Both kinds of thermal
expansion are characterized by the average coefficients of linear
ai and cubic, or volume, av expansion for the given temperature
range.'.0.2] THERMAL EXPANSION 581
40.2.2 If l0 is the length of a body at the temperature 0 °C, its
(donga t ion AI upon being heated to the temperature t °G is
M = a tl0t.
from which a; = Alll0t.
The linear expansion coefficient characterizes the elongation per
unit length Al/l0 as the body is heated one degree. For most solids
aL  10-6 to 105 K1 and depends only slightly on the temperature.
40.2.3 When a body is heated from 0 °G to t °C, its volume
increases from V0 to V according to the law
V = V0 (1 + Oyt),
from which av — V/V0t. The cubic expansion coefficient characterizes
the volume increase per unit volume A VlV0 as the body
is heated by one degree.
In a first approximation, the relation between the coefficients av
and at is of the form
av  3a;. /
40.2.4 The thermal expansion of solids is due t<Ahe anharmonici-
l y of the thermal vibration of the particles in me lattice (40.1.4).
I f r0 is the equilibrium distance between adjacent particles, then,
at some arbitrary instant, the distance betweenxjdiem is r — r0 +
|- x, where x is’ the mutual displacement of thhparticles from
tlieir equilibrium positions due to thermal vibration. The forces
between the particles in the lattice are not quasi-elastic, but
depend on the displacement x according to the relation
Fx = —kx + bx2,
where x is the coefficient of quasi-elastic force (40.3.5), and b
is the coefficient of anharmonicity of the vibration. The term bx2
characterizes the deviation of the vibration from its harmonic
type due to the different dependences of the forces of attraction
and repulsion on the distance.
When a solid is in an equilibrium state the positions of the crystal
lattice points should not vary with time and the average force
acting on each particle should equal zero:582 STRUCTURE AND PROPERTIES OF SOLIDS [40
If the vibration of the particles was strictly harmonic (Fx = —ycx),
the average displacement of the particles would be (x) =
= — (Fx)/k = 0, i.e. there would be no thermal expansion. For
real anharmonic vibration, it follows from the condition that
(Fx) = 0, that
— K (x)b (x2) = 0 or (a;) = — (x2).
x
For thermal vibration with small amplitudes the potential energy
Ep of the vibrating particles is equal (28.2.3) to:
But, according to the law of equipartition of energy (10.4.6),
Ev = kTl2, where h is Boltzmann’s constant (8.4.5) and T is
the absolute temperature. Thus
x (x2)
kT
~2 ~ ~ or
Finally, (x) = bkT/yi2. Hence, the average distance between
the particles of a solid increases when it is heated, leading to
thermal expansion.

\section{Brief Information on the Elastic Properties of Solids}
\label{sec-40.3} 


40.3.1 Deformation of a solid is the change in its dimensions and
volume. Deformation is usually accompanied by a change in the
shape of the body. Sometimes (in triaxial tension or compression)
no change in shape occurs. Deformation may be caused by external
forces acting on the body, a change in its temperature (40.2.1)
and other reasons.
Deformation of a body leads to displacement of its particles from
their initial equilibrium positions at the crystal lattice points
(8.1.1) to new positions. This displacement is opposed by the
forces of interaction between the particles. Internal elastic forces
are developed inside the body to counterbalance the external
forces causing the deformation.40.3] ELASTIC PROPERTIES OF SOLIDS 583
40.3.2 Deformation is said to be elastic (elastic deformation) if it
disappears as soon as the external forces causing it cease to act.
At this the particles of the solid return to their initial positions
of equilibrium (40.3.1). Inelastic deformation of a solid leads to
irreversible rearrangement of the crystal lattice (40.1.1) and the
shape of the body is not restored when the external forces are removed.
Such deformation is said to be plastic (40.3.8). The transformation
of elastic deformation into plastic deformation may occur
upon prolonged action of even small external forces applied to the
body. The reverse transformation cannot occur.
The degree of elasticity is the ratio of the work that can be done by
a deformed body when the deforming forces are slowly removed
to the work done by external forces in deforming the body.
40.3.3 The physical quantity equal to the elastic force dFei per
unit area dS of the cross section of the body being deformed is
called the stress a:
The stress is said to be normal if force dFei is normal to the surface
dS and tangential, pr shearing, if the force is tangential to this surface. 
40.3.4 A measure of the deformation is the unu strain, or relative
deformation. It is equal to the ratio of the absolute deformation
Ax to the initial value of the quantity x, which characterizes the
shape or dimensions of a deformed body: Axlx.
Hooke's law states that the stress in an elastically deformed body
is proportional to its unit strain. Thus
c = K: Ax
where Kx is the modulus of elasticity, equal in magnitude to the
stress developed in relative deformation equal to unity. The
quantity ax = 1 lKx is called the coefficient of elasticity. Hooke’s
law is valid only for sufficiently small relative deformation. The
stress crp at which the deformation is no longer proportional to the
stress is called the proportional elastic lim it (point A in Fig. 40.1).
40.3.5 In addition to elastic forces, there are forces of other origin
than elastic forces, but which satisfy the equation
Fx = —Kx,584 STRUCTURE AND PROPERTIES OF SOLIDS [40
where Fx is the projection of the force on the direction in which
the absolute linear deformation x takes place. They are called
quasi-elastic forces. The quantity x is called the coefficient of the
quasi-elastic force.
40.3.6 The simplest type of deformation is longitudinal, or uniaxial,
extension (or compression), which consists in elongating
(shortening) a body by the action of the tensile (compressive)
force F. Elastic extension (compression) ceases when F = Fe,
where Fe] is the elastic force (40.3.1). The unit strain (relative
deformation) A xlx = A III, where AI is the change in length due
to the action of force F, and I
is the initial length of the body.
C
According to Hooke’s law
where the modulus of elasticity
Kx = E is called Young's modulus.
Young’s modulus is equal to
the normal stress a at which the
linear dimension of the body is
Fig. 40.1
doubled: AI = I*.
40.3.7 The dependence of the
normal stress a on the unit strain
A Hi in uniaxial extension is called the tensile stress-strain diagram
(Fig. 40.1). Beyond the proportional elastic limit (40.3.4), an
increase in a leads to a considerable increase in A III. At the stress
cry, corresponding to point in the'stress-strain diagram, the unit
strain continues to increase without increasing the stress (horizontal
portion BB' in the diagram). The stress ay is called the
yield point. The maximum stress at that the body can withstand
before rupture, point C in the diagram, is called the ultimate
strength or, if the body is in tension, the tensile strength. Point D
corresponds to the rupture of the body.
40.3.8 If the load on a body, deformed to the stress oa (point a
in the diagram in Fig. 40.1) is slowly removed, the curve a =
— cp (AZ/Z) is the straight line aB, parallel to the straight portion
OA of the diagram. The line segment OR determines the permanent
et, typical of plastic deformation.
s
* Assuming that Hooke’s law is valid for such large deformation.'.u/i] PHASE TRANSITIONS IN SOLIDS 585
40.3.9 Shear is the deformation of a body in which all its flat
layers, parallel to a certain shear plane, are displaced parallel
In one another without distortion or any changes in dimensions
(Fig. 40.2). Shear results from
I he action of tangential force F,
applied to face BC, parallel to
I he shear plane. Face AD, parallel
to face BC, is fixed. For
small shear deformation
CC
y  tan y =
CD 77777777777777777777'/
where CC' = Ax is the absolute
shear, y is the shear angle, or Fig. 40.2
unit shear strain, expressed in rail
ians. According to Hooke’s law
40.3.4), the unit shear strain is proportional To the tan/gentiall
(shearing) stress aT = F/S, where S is the area of woJZC, i.e.
oT = Gy.
The quantity G is cafled the shear modulus. Its magnitude is equal
to the tangential, or shear, stress that produces a unit shear strain
equal to unity in the specimen.

\section{Basic Concepts of Phase Transitions in Solids}
\label{sec-40.4}

40.4.1 The heating of a solid crystalline body leads to an increase
in the amplitude of the anharmonic thermal vibration at the points
of the crystal lattice (40.2.4) and to an increase in the average
interpoint distances in the lattice (thermal expansion, Sect.
40.2.1). Intense heating destroys the crystal lattice and leads
to the transition of the substance from its solid to its liquid or
vapour phase (first-order phase transition, Sect. 12.4.2).
40.4.2 The transition of a substance from the solid to the gaseous
state is called sublimation. Transition from the solid to the liquid
state is called melting, or fusion. Melting begins at a definite
temperature Tm for the given pressure; it is called the melting
point. This temperature remains constant during the melting586 STRUCTURE AND PROPERTIES OF SOLIDS [40
process. Shown in Fig. 40.3 is the dependence T ((?), where Q is
the quantity of heat (9.2.1) delivered to the solid being heated.
The isothermal portion BC corresponds to a two-phase system:
solid-liquid (cf. Sect. 12.3.2, where the two-phase system consists
of a liquid and saturated vapour). In the melting process, the
substance is converted from a more ordered crystalline state to a
less ordered liquid state. According to the second law of thermodynamics
(11.3.2), melting is associated with an increase in the
entropy (11.4.3) of the system.
40.4.3 The quantity of heat required to melt unit mass of
a solid at the temperature Tm is called the heat of fusion, or of
melting. It is also called the latent heat. It follows from the first
law of thermodynamics (9.3.1) that
rm = u — us + P (*>1 — i>s)>
where u and us are the specific internal energies (9.1.2) of the
substance (internal energies per unit volume) in the liquid and
solid phases, and vs are the specific volumes of the liquid and
solid phases, and p is the constant pressure at which the first-
order transition proceeds (12.4.2).
40.4.4 The melting point Tm depends upon the pressure p.
A change of dp in the pressure leads to a change in the melting'<0.4] PHASE TRANSITIONS IN SOLIDS 587
point by dTm. The dependence Tm = Tm ( p) is expressed by the
Clapeyron-Clausius equation (see also Sect. 13.6.5):
dTm (v — ys)
dp rm
The notation is defined in Sect. 40.4.3. The Clapeyron-Clausius
(equation is valid for any first-order phase transition.
In the overwhelming majority of cases the melting of a substance
leads to an increase in its specific volume, i.e. v > vs and
(lTm/dp > 0 (rm >0), i.e. the melting point increases with the pressure. Such substances can be in the solid state at high pressures
in the temperature region above the critical temperature
(12.3.1) .
The specific volume of certain substances (water, bismuth, gallium)
increases in melting, so that v < vs. Ice, for instance,
at 0 °C has a lower density
than water. For such substances
dTmldp < 0, i.e. the melting
point decreases as the pressure
is increased.
40.4.5 In a pT diagram, the dependence
of the temperature on
the pressure p = f (T) in a first-order
phase transition is represented
by a curve on which each point
corresponds to equilibrium of
two existing phases. The vapour
generation curve, for example,
0 T
shows the equilibrium of the
two-phase system: liquid-vapour.
Fig. 40.4Since vY — v > 0 (where vy is
the specific volume of the vapour),
it follows from the Clapeyron-Clausius equation (40.4.4), that
dp/dT > 0. The curve p = f (T), separating the regions of the
liquid and gaseous states of a substance, ends at critical point C
(12.3.1) (Fig. 40.4). Shown in Figs. 40.5 and 40.6 are the equilibrium
curves of the two-phase solid-liquid states for the cases
when v] > vs and u < vs, respectively.
40.4.6 In cooling liquids to a certain temperature, called the
freezing point T (or temperature of solidification), the substance is
solidified (crystallized, or frozen). This is the transition from588 QUANTUM PHYSICS OF SOLIDS [M
a liquid to a solid crystalline state. In this process a quantity of
heat is evolved that is equal to the heat of melting (40.4.3). The
freezing point is equal to the melting point and depends upon the
presence of impurities in the substance. Impurities lower Tf if
they do not form mixed crystals with the substance. These are
called solid solutions. Sea water, for instance, containing dissolved
ip Vi >vs
P Vt <Vs
Liquid
Solid Liquid
Solid
0 T 0 T
Fig. 40.5 Fig. 40.6
salts, freezes at a lower temperature than distilled water. This
serves as the basis for producing cooling mixtures.
40.4.7 The freezing (crystallization) of liquids is associated with
a change in the nature of the thermal motion of their particles.
The time the particles remain in a temporary settled state (relaxation
time, see Sect. 13.2.4) is increased. The forces of mutual
attraction between the particles change the thermal motion into
random thermal vibration about the points of the crystal lattice. The transition of the substance to a more ordered phase is associated
with a reduction in the entropy of the system (11.4.3).

% Local Variables:
% TeX-master: "handbook-physics"
% coding: utf-8
% End:
